\section{全幾何学的H3の判定と表示の四成分}
\label{sec:cubic-presentation-components}

\begin{theorem}\label{thm:h3-complete-geometric-admissibility}
標数零の代数閉体上で、強い内在的H3の全対象は以下の表示から得られる。
各行の判定は、表示された中央又は正規三次直線がH3切断を与えるための
必要十分条件である。
\begin{enumerate}
\item 符号重複度 $(1,3)$ は定理\ref{thm:h3-order2-full-arrows}の全係数族、
$(2,2)$ は定理\ref{thm:h3-order2-22-L-complete}で認識した
曲がったClifford族の許容開である。
\item 一交換子CYは定理\ref{thm:cubic-CY-one-commutator-presentation-closure}の
三乗根Ore表示に限られる。二元二次補正の階数 $0,1,2$ にはそれぞれ
定理\ref{thm:all-skew-H3-geometric}、
\ref{thm:cubic-CY-one-segre-square-split-resonance-exact}、
\ref{thm:cubic-CY-one-segre-double-split-outer-H3}を適用する。
\item 可換包絡は滑らかな可換三次曲面を与える。
三交換子CYの対称 $\alpha$ は系
\ref{cor:cubic-CY-three-alpha-symmetric-complete}の許容表示に限られる。
\item 三交換子CYの非対称 $\alpha$ では、
尖点・接線・共点型は定理\ref{thm:cubic-CY-three-alpha-additive-H3}のKleinian条件、
節点型は定理\ref{thm:cubic-CY-three-alpha-nodal-E6-root-discriminant}の
$E_6$ 根判別式で判定する。
三角形の体積切断と割線では、位数によらず定理
\ref{thm:cubic-CY-three-alpha-nontorsion-full-folded}の平方自由条件に従う。
原始三乗根での三角形の全追加中心方向は
定理\ref{thm:cubic-CY-three-alpha-primitive-triangle-test}で判定する。
滑らかな楕円境界は、有限位数では定理\ref{thm:finite-order-elliptic-admissibility}の
中心と単純表現の判定、無限位数では定理
\ref{thm:cubic-CY-three-alpha-elliptic-diagonal-all-nontorsion}の八次判別式に従う。
\end{enumerate}
全矢は、生成空間の線型同型で二次関係部分空間と三次切断直線を移すものを
共通次数スカラーで割って得る。
従って同一表示に対するこの安定化群が内在的慣性を与える。
\end{theorem}
\begin{proof}
三次関係の必要性と二次包絡の関手的復元は
命題\ref{prop:intrinsic-ring-isomorphisms}と正規三次incidenceによる。
非CY包絡の全スペクトルと全Jordan型は次定理の被覆の通りである。
CYの場合、零と二交換子は既に空で、
四交換子以上は可換多項式包絡に限られる。
一交換子の包絡は前掲の三乗根Ore表示で尽くされる。
三交換子CYの非 $\alpha$ 部分の排除は系
\ref{cor:cubic-CY-three-presentation-closure}の証明に列挙した。
残る $\alpha$ では、対称の場合は前掲の全許容判定、
非対称の場合は滑らかな三次、節点、割線、三角形、尖点、接線、共点の被覆を使う。
零先頭および重複直線、成分を巡回する非許容三角形は
各々の排除命題で消えている。
各行で挙げた定理が必要性と十分性を与えるので、
包絡と切断係数の全幾何学的点を判定できる。
矢の記述は同命題により二次包絡・切断の同型と一致する。
\end{proof}

\begin{corollary}\label{cor:h3-rank-from-classified-smooth-domain}
標数零の代数閉体上で、$B$ を左右noetherian quadratic Koszul AS4型14641とし、
$G$ を $Gv=\mu_B^{-1}(v)G$ を満たす非零三次元とする。
$R=B/(G)$ が整域で、$\qgr R$ の大域次元が有限なら、
$\pi R$ はGabriel--Krull rank-one torsion-freeであり、$(\qgr R,\pi R)$ はH3である。
\end{corollary}
\begin{proof}
$B$ は整域なので $G$ は正則正規。商はAS Gorenstein次元3、parameter 1、
$\mu_R=1$ である。包絡による有限自由解消と飽和の比較から
$\chi$、sectionの復元、例外性が従う。
有限大域次元の仮定でAS Serre双対は全有界coherent導来圏上の
$(-1)[2]$ となるので、大域次元はちょうど2である。
従ってここで残っているのは二段Gabriel--Krull商によるrankである。

前定理の幾何学的被覆では、このrankを排除の入力に使う箇所を次のように置き換えられる。
まず中山作用の四・三固有値、単冪部分および二固有値の除外に使うのは
Frobenius対合、Hilbert階数、正規線型式、および整域性である。
三次元の自己Extに関わる除外では
補題\ref{lem:cubic-odd-Serre-determinant}のSerre標的と
そのshift作用を用いる。ここでも上で得た双対が同じである。
これにより非CYは $(1,3),(2,2)$ の二Clifford表示へ入る。
これらの幾何判定では、一般中心での斜体条件に整域性を使い、
残る局所条件に有限大域次元を使う。
従って各許容表示に対して既に証明したrankが従う。

CY部分では、可換な場合を除き、零と二交換子の排除および一交換子の
原始三乗根模型への還元は、中央三次の線型核と整域性、
または同じ点のSerre障害を用いる。
Sklyaninと指定cocycle twistのrank障害は具体的な三次商の非整域性であり
（Sklyaninと標準 $V_4$ の排除）、ここでも排除される。
三交換子で $\alpha$ 以外が消えるのは
系\ref{cor:cubic-CY-three-presentation-closure}の冒頭に列挙した
二次tensor階数と中央三次の排除による。
残る $\alpha$ では対称・非対称それぞれの表の必要被覆と
有限大域次元の判定を適用する。アフィン環の大域次元が必要な場合は
中央一次による局所化が高次Ext消滅を保つことを使う。
したがって滑らかな楕円、節点、割線、三角形、尖点・接線・共点、
および対称Cliffordの各許容模型のどれかに入る。
各模型では、射影中心上の一般斜体、または境界での無限降鎖から
Gabriel--Krull rankを別に証明している。
可換包絡のsmooth切断でもrankは通常の整な曲面のrankである。
これで任意の表示に対して所要のrankが従う。
\end{proof}

\begin{theorem}\label{thm:h3-presentation-four-components}
四次元生成空間上の六二次関係と切断直線によるincidenceに、H3の点の部分位相を入れる。
基底変更を法とした幾何学的表示の位相には、開かつ閉な既約成分がちょうど四つある。
一般幾何商の次元と成分を特徴づける条件は次の通りである。
\[
\begin{array}{c|l|c}
\operatorname{tr}(\mu_B|_{B_1})&\text{条件}&\text{次元}\\\hline
4&\text{一交換子CY}&3\\
4&\text{三交換子CYおよび可換三次曲面}&8\\
-2&\text{符号の重複度 }(1,3)&4\\
0&\text{符号の重複度 }(2,2)&5
\end{array}
\]
ここで交換子の数は、二次関係空間と $\bigwedge^2 B_1$ の交叉次元である。
四成分の間にH3表示を保つ特殊化や交叉はない。
\end{theorem}
\begin{proof}
固定四次元空間 $V$ と関係部分空間 $I_2\subset V^{\otimes 2}$ に対し、
二次双対 $E=T(V^*)/(I_2^\perp)$ を取る。
次数3と4の関係挿入行列のrankを固定する局所閉stratumでは、
$E_3,E_4$ は有限局所自由で、その形成は底変更に両立する。
全H3表示が属するrank stratumでは、それぞれの階数は4と1である。
Frobenius対合の開条件の下で、二つの積
\[
 V^*\otimes E_3\longrightarrow E_4,\qquad
 E_3\otimes V^*\longrightarrow E_4
\]
はともに完全である。従って $xy$ と $y\nu(x)$ の $E_4$ での同一視で
定まる一次作用 $\nu$ は、両対合の行列の積と逆行列で正則に与えられる。
$E_4$ の自明化の変更は二つの行列に同じ倍率をかけるため、作用は貼り合わさる。
全H3表示はKoszul AS4なのでFrobenius条件を満たし、
$\mu_B|_V=-\nu^*$ である（\cite[Theorem 6.3の証明]{RRZ}）。従って中山traceは
この有限次数stratumから制限した正則関数で、基底変更で不変である。

定理\ref{thm:cubic-three-eigenvalue-sector-empty}、
\ref{thm:cubic-four-simple-sector-empty}、
\ref{thm:cubic-jordan22-unipotent-empty}、
\ref{thm:h3-order3-two-eigenvalues-excluded}により、
残るのはCYおよび符号スペクトルである。
符号重複度$(3,1)$の全候補は
定理\ref{thm:h3-positive3-gamma-zero-exclusion}、
\ref{thm:h3-positive3-canonical-base-coverage}、
\ref{thm:h3-positive3-cone-base-exclusion}、
\ref{thm:h3-positive3-Jordan-base-exclusion}、
\ref{thm:h3-positive3-diagonal-base-exclusion}で排除される。
従ってtraceの値は $4,-2,0$ に限られ、各逆像は相対的に開かつ閉である。

trace $4$ の二成分は系\ref{cor:cubic-CY-presentation-two-components}である。
trace $-2$ は定理\ref{thm:h3-order2-full-arrows}の既約四次元表示商で尽くされる。
trace $0$ では系\ref{cor:h3-order2-22-rank2-all-Jordan-coverage}が
混合二次関係数2を排除し、定理\ref{thm:h3-order2-22-L-complete}が残る全表示を
曲がったClifford族へ認識する。
定理\ref{thm:h3-order2-22-clifford-arrows}の既約開係数空間とその基底変更の像は
この全部分を覆い、商は一般次元5となる。
いずれも既出の族の表示をincidenceに代入する写像であり、その像は連続である。
よって四部分は全H3幾何表示の既約成分である。
\end{proof}

この位相は二次包絡と三次直線を持つ表示から定めたものである。
定義\ref{def:h3-intrinsic-family}の相対内在族との同値および有効降下は
定理\ref{thm:h3-intrinsic-algebraic-stack}で扱う。
この相対定義には二次包絡の三次平坦性が含まれる。
より弱い平坦Abel圏変形からこの条件が自動的に従う範囲は、別の問題である。

\subsection{次数スカラー除去と表示の有効降下}
\begin{proposition}\label{prop:h3-presentation-canonical-degree-descent}
全次数平坦な包絡と切断を持つ任意の基底変更不変な表示部分を取る。
この表示を $\mathrm{PGL}_4$ で割ったとき、次数付き環を降下させるための
$\mathbb G_m$-gerbeは中立である。
具体的に、一次束 $V$ と切断直線 $\mathcal G$ に対して
\[
 \mathcal L=\det(V)\otimes\mathcal G^\vee,\qquad
 B^{\rm can}_n=B_n\otimes\mathcal L^{-n},\qquad
 R^{\rm can}_n=R_n\otimes\mathcal L^{-n}
\]
と置けば、包絡・切断・商は $\mathrm{PGL}_4$ 商へ有効に降下する。
この標準化は非被約底を含む任意の基底変更に両立する。
\end{proposition}
\begin{proof}
枠付き表示の空間では、$cI$ は $B_n,R_n$ に $c^n$、
切断直線に $c^3$、$\det(V)$ に $c^4$ で作用する。
従って $\mathcal L$ 上の作用は $c$ であり、
$B^{\rm can}_n,R^{\rm can}_n$ の線型化は中心スカラーを通じて自明となる。
積は $\mathcal L^{-m}\otimes\mathcal L^{-n}=\mathcal L^{-(m+n)}$ により
同じように降下する。二次関係部分束、三次切断直線、商写像も次数に従って
同じ線型化を持つので、全表示を降下できる。
別の言い方では、元の $\mathrm{GL}_4$ 商のgerbeに重み一の可逆層
$\mathcal L$ があり、ここで与えた捩りはその重みを打ち消す中立化である。
降下した一次束と切断直線には
$\det(V^{\rm can})\otimes(\mathcal G^{\rm can})^\vee\simeq\mathcal O$
が備わる。構成には線型化と可逆層のテンソルだけを使うため、
基底変換の平坦性や被約性を仮定しない。
\end{proof}

\begin{proposition}\label{prop:h3-presented-qgr-fpqc}
$R$ を可換noetherian環 $A$ 上の次数付き代数とし、
調べる全noetherian基底変更上で右noetherianとする。
このとき $\qgr R$ の対象と射には、affineな忠実平坦基底変更に関して
有効降下が成り立つ。特に上の標準化で降下した表示の
構造対象 $\pi R$ とshiftを含むtails圏は、表示のfpqc降下と両立する。
\end{proposition}
\begin{proof}
全次数付き加群の商 $\operatorname{QGr}R$ と、その飽和section関手 $\omega$ を使う。
加群 $M$ のtorsion部分 $\tau M$ は各次数で
$\varinjlim_n\operatorname{Hom}_{\mathrm{Gr}R}(R/R_{\ge n},M(\ell))$ で計算できる。
$\omega\pi M$ は同じ次数で
$\varinjlim_n\operatorname{Hom}_{\mathrm{Gr}R}(R_{\ge n},(M/\tau M)(\ell))$
で計算できる。
第一引数はいずれも有限表示なので、平坦底変更は両式と可換する。
次数付き加群のfpqc降下により、飽和section加群の降下は有効である。
降下した次数付き加群が飽和していることは、飽和への単位射が
忠実平坦化後に同型であることから従う。
よって飽和加群の充満部分圏である $\operatorname{QGr}R$ の対象・射も降下する。

降下した飽和加群を $M$ とする。忠実平坦化後の対象がcoherentなら、
そこで有限個の斉次元が生成する部分加群から $M$ への射が
tails商で全射になる。それら有限個の元に現れる元の $M$ の斉次成分で
有限生成部分加群 $N$ を作れば、$M/N$ は忠実平坦化後にtorsionとなる。
各元に対する十分高次の消滅は忠実平坦性で反映されるので、$M/N$ 自体がtorsionである。
従って $\pi M=\pi N$ はcoherentである。
構造対象とshiftは自由次数付き加群と次数移動から得られ、同じ降下に両立する。
\end{proof}

\subsection{全表示の開性と有限型表示stack}
\begin{proposition}\label{prop:h3-primitive-triangle-constructible}
原始三乗根三角形の六包絡係数と全切断直線の有限型係数空間では、
定理\ref{thm:cubic-CY-three-alpha-primitive-triangle-test}の
幾何学的H3条件は構成可能である。
\end{proposition}
\begin{proof}
構成可能性は、各被約既約な閉部分上で幾何的一般点と同じ判定を
稠密開集合へ広げ、noetherian帰納法を行うことで示す。
以下、幾何的一般点上の有限個のデータが必要なら、
有限体拡大上に定義して有限型の優越する底へ移る。
その上の稠密開像には元の底の稠密開集合が含まれるので、
幾何学的条件についてこの変更で十分である。

中心商 $C$ の幾何整域性は有限表示族で構成可能である。
$C$ が整域である範囲で、有限 $C$ 加群 $U$ のtorsionを次のように広げる。
一般点上の分数体基底を選び、分母を払って
$U\to C^r$ を作る。その核はある非零元で殺される。
写像、核、像、余核を底上に広げ、generic flatnessの開上では
この有限加群の核と写像が底変更と可換する。
分母を消す元が各fiberで非零である開に制限すると、
この核は各fiberのtorsionそのものである。
核の台の底への像は構成可能なので、一般点における零・非零の
いずれも稠密開へ広がる。
同様に一般中心上の次元、正則trace対合、中心次元の階数は、
分母を払った有限行列の小行列式が各fiberの一般点で
零かどうかで判定する。非零小行列式はその非消滅集合の像を使い、
恒等的に零のものは一般点での恒等式を底に広げればよい。

九次元中心単純である範囲では、そのBrauer類の非零性を調べる。
幾何的一般点で分裂するなら、行列環への同型を有理関数体上で選び、
有限個の分母、積の恒等式および可逆行列の条件を広げる。
従って近傍の幾何fiberも分裂する。
幾何的一般点で非分裂なら、定理
\ref{thm:cubic-CY-three-alpha-primitive-triangle-test}の証明により
滑らかな射影モデル上で非零剰余を持つ素因子が一つある。
そのモデル、素因子の滑らかな稠密開、巡回symbol表示および
剰余を記述する単元を広げる。
モデルは稠密開上で滑らかな射影族となり、
関数体を同定する二つの稠密アフィン開を同時に広げれば、
各幾何fiberでも元の中心の射影モデルとなる。
素因子の一般点で局所方程式 $z$ を選び、symbolの各成分を
$z$ の整数冪と単元の積に書く恒等式も広げる。
剰余が零でないことは、その剰余単元の三乗根を加えた
有限\'{e}tale曲線被覆が幾何的に整であることに等しい。
有限表示族の幾何整域性の構成可能性により、この非零剰余が稠密開に広がる。
よって分裂・非分裂の区別はnoetherian帰納法で構成可能である。

射影四chartの判定には、各 $A_f$ の次数別構成から得る有限中心表示を使う。
$A_f$ は、平坦な斉次商の切断成分を中央元で局所化して得るため、
係数底上で平坦であり、底変更に両立する。
その包絡環 $A_f^{\rm op}\otimes A_f$ は可換有限型環上有限なのでnoetherian。
対角の有限自由解消を必要な第四項まで取り、
各核の底上平坦性から全幾何fiberで同じ解消を使用できる。
従って単純表示の次数を $1$ から $27$ まで動かす有限個の式、
単純性の階数と第三Extの階数を一つの係数族上で課せる。
射影による構成可能像の補集合が第三条件を与える。
\end{proof}

\begin{theorem}\label{thm:h3-presentation-algebraic-open}
定理\ref{thm:h3-finite-degree-normal-incidence}の有限型schemeを $X$ とする。
全H3幾何表示を与える部分は $\mathrm{PGL}_4$ 不変なZariski開
$X^{\mathrm{H3}}\subset X$ である。
従って、幾何fiberがH3である全正規三次表示は
有限型代数stack
\[
 [X^{\mathrm{H3}}/\mathrm{PGL}_4]
\]
をなす。これは任意底上の次数付き表示とその基底変更の商であり、
各幾何点で、系\ref{cor:h3-artin-intrinsic-formal-atlas}の仮定を満たし、
二次包絡の三次成分が平坦である内在Artin変形の群oidを与える。
その位相は定理\ref{thm:h3-presentation-four-components}の四成分を持つ。
\end{theorem}
\begin{proof}
まず幾何学的H3集合が構成可能であることを示す。
定理\ref{thm:h3-complete-geometric-admissibility}の係数模型を用いる。
二つの非CY模型は、既出の射影conic判定で定まる許容開である。
一交換子模型は二元二次補正の階数で三stratumに分け、
既出の純立方・二元重根・局所orderの有限判定を使う。
可換三次では普通の三次formの滑らかな開、
対称Clifford模型では中心のJacobianと少数階数のHessian条件である。
この条件の失敗は有限個の中心座標・階数条件を持つincidenceの射影である。

非対称 $\alpha$ の滑らかな楕円部分は
定理\ref{thm:h3-elliptic-mixed-order-projective-open}の開。
量子三角形の体積切断・割線は
定理\ref{thm:cubic-CY-three-alpha-nontorsion-full-folded}の判別式開である。
原始三乗根の追加切断は命題\ref{prop:h3-primitive-triangle-constructible}で扱った。
節点では $E_6$ の積1スペクトルのトーラスからの有限係数式と、
有限個の根指標の非消滅を使う。
尖点・接線・共点では、対応する三型それぞれのPBW量子化との
filtered同型と根条件を表す有限式
（命題\ref{prop:cubic-CY-three-alpha-additive-affine-test}の証明）を使う。
この三型の包絡のAS正則性には、固定出典\cite{AS4CYpinned}の
各型の全係数PBWの十分性を用いる。
いずれも有限型係数scheme上の有限個の構成可能条件、またはその像である。

各係数模型から六二次関係と切断直線への写像は、
表示した係数を有限次数の関係行列へ代入して作る。
階数を実現する小行列式で有限個の開に分けると、
関係部分空間と切断直線が束となり、$X$ への写像が得られる。
必要なら模型の構成可能集合を有限個の局所閉部分に分けて
被約構造を取れば、所定階数を超える小行列式は恒等的に零になる。
$\mathrm{PGL}_4$ の基底変更も同時に動かしてChevalleyの定理を適用すると、
全被覆は $X$ の構成可能集合となる。

次に一般化で安定であることを示す。
$y$ が $x$ の一般化で $x$ がH3なら、
$y$ の被約閉包を $x$ で局所化し、極大イデアルのblow-upを正規化する。
$x\ne y$ ならその例外支持の余次元1点の局所環は離散付値環である。
このtraitは $x$ に特殊化し、一般点は $y$ に行く。
係数schemeは標数零体上有限型なので正規化は有限である。
幾何的H3判定は代数閉体拡大で不変だから、特殊fiberへ適用できる。
系\ref{cor:h3-trait-homological-generization}により、その幾何一般fiberもH3。
従って構成可能集合は一般化で安定し、Zariski開である。

全H3fiberを持つ任意底上の表示では
定理\ref{thm:h3-finite-degree-normal-incidence}により全斉次成分が有限局所自由。
一次枠と共通次数スカラーを法とする商は上の代数stackである。
重み一のorientationによる降下で表示は任意底変更と可換し、
内在Artin変形との同定は
系\ref{cor:h3-artin-intrinsic-formal-atlas}による。
最後に開集合は従来の幾何表示の部分位相を持ち、
基底変更商の位相へ四成分の計算が移る。
\end{proof}

\begin{theorem}\label{thm:h3-presentation-strong-noetherian}
開 $X^{\mathrm{H3}}$ とその代数stackから得られる全包絡・斉次商は、
任意のnoetherian底変更上で左右noetherianである。
従って構造対象付きcoherent tailsの族には、noetherian底上の
affine忠実平坦写像に関する有効降下が成り立つ。
\end{theorem}
\begin{proof}
最初に、定理\ref{thm:h3-complete-geometric-admissibility}の模型は
noetherian係数底上の強noetherian環として使用できることを確認する。
ここで必要なのは、各模型の係数条件を満たす係数底への引き戻しが
どれもnoetherianになることである。
\begin{itemize}
\item 可換模型は多項式環である。一交換子模型は前掲の逐次Ore提示を持つ。
その自己同型の係数は単元なので、任意のnoetherian底上でOreのHilbert基底定理が適用できる。
\item 符号二系統と対称Cliffordの包絡では、各平方、または表示した
有限個のその線型結合が中央である。前掲の有限中心基底は
同じ順序語の還元から任意底上の有限生成を与える。
原始三乗根三角形でも $t,X,Y,Z$ の上で同じ順序語の還元を使える。
従って各環は有限型の可換環上有限であり、全てのnoetherian底変更でnoetherianである。
\item 尖点・接線・共点の三包絡は前掲のPBW順序基底を持ち、
指定した重みで随伴環が可換多項式環となる。
活性三変数の正重みと中央変数の零重みの下で、
最高記号からの帰納法により任意底上でnoetherianである。
\item 残りの非対称 $\alpha$ 模型では、中央正則な一次元 $t$ の商 $S$ と、
$S$ の中央正則な三次元 $g$ の商 $H$ を順に使う。
滑らかな楕円、量子三角形、量子割線および節点の各 $H$ は、
前掲の点行列のグラフが与える捩れ斉次座標環である。
この同定は係数底上でも成り立つ。
実際、点行列の階数2の核と逆核は相対的に作れ、
点曲線は平面三次の相対Cartier因子である。
各fiberが被約Gorenstein三次曲線で、正次数の捩れ線束は全成分に正の次数を持つため、
その $H^1$ は双対性から零となる。
従って正次数sectionは底変更に両立する。
被約の普遍係数底上では、関係挿入行列の各次数の一定rankと
正則中央元の倍写像から $S,H$ は次数別有限局所自由である。
従って生成元から捩れsection環への各次数の写像は
fiber上の同型から相対的にも同型で、任意底変更する。
$H$ の線束列は相対ample filterである。底をaffineとして、
任意の固定線束$\mathcal F$に対し、
$\mathcal F\otimes\mathcal O_E(m)$が相対的に大域生成される$m$を選ぶ。
幾何fiberの既約成分$C$は平面三次の成分だから
\[
 \deg(\mathcal F|_C)\ge -m\deg(\mathcal O_E(1)|_C)\ge-3m.
\]
各捩れ$\sigma^{j*}\mathcal O_E(1)$の成分次数は少なくとも1なので、
長さ$n$の捩れ積$\mathcal L_n$に対し
$\deg((\mathcal F\otimes\mathcal L_n)|_C)\ge n-3m$となる。
従って$n>3m$で全幾何fiberの全成分の次数が一様に正となり、
相対的にampleである。affine底上ではampleとなる。
従って \cite[定理1.3, 7.14]{KeelerAmpleFilters}により
任意noetherian底上の $H$ は左右noetherian。
正次数の中央正則元 $g,t$ の順の持ち上げで、
$S,B$ も同じである。
\end{itemize}
いずれの場合も斉次商は包絡の二側イデアルによる商なので、同じ性質を持つ。
係数空間を構成可能stratumに分けて定義した表示では、対応する被約係数stratumを用いてよい。

次に、この模型による幾何被覆をnoetherian底上へ移す。
$A$ を任意のnoetherian環とし、$T$ を問題の包絡または斉次商とする。
$T$ は全次数で有限局所自由である。
まず $A$ が整域であるとする。
その分数体の代数閉包上で、表示は上記の模型のどれかに同型である。
模型の係数、基底変換およびその逆は有限個なので、
ある有限体拡大上に定義できる。
その分母を払えば、適当な濃密主開上の有限忠実平坦拡大上で
表示が模型の引き戻しであることを表す有限係数恒等式を得る。
ここで有限体拡大の生成元を整にし、有限格子を一度局所化して自由にできる。
係数条件が局所閉である場合は、非零条件も同じ局所化に含める。
従ってこの主開で $T[f^{-1}]$ は忠実平坦降下により左右noetherianである。

閉補集合へnoetherian帰納法を適用して $T/fT$ をnoetherianとする。
$A$ の $f$ 進完備化 $\widehat A_f$ は平坦であり、
$T\otimes_A\widehat A_f$ の $f$ 進随伴環は $(T/fT)[s]$ である。
正則な $f$ を用い、命題
\ref{prop:h3-noetherian-local-lift}の証明と同じ次数別の有限自由性、
完備分離性、初期式の逐次消去により、完備化後の環は左右noetherianとなる。
$A[f^{-1}]$ と $\widehat A_f$ の二つの像は $\operatorname{Spec}A$ を覆うので、
その積からの忠実平坦降下で $T$ は左右noetherianである。

被約底では、有限個の極小素イデアル $\mathfrak p_i$ に対し
$A\to\prod_i A/\mathfrak p_i$ は単射である。
次数別平坦性で $T$ は各 $T/\mathfrak p_iT$ の有限直和に
右加群として埋め込まれる。各直和因子は前段によりnoetherian右加群なので、
$T$ 自身も右noetherian、左も同様である。
非被約底のnilradical $\mathfrak n$ は冪零であり、
その有限濾過の各商は $T/\mathfrak nT$ 上有限生成である。
従って被約底の場合から同じ結論が従う。
この議論は任意のnoetherian底変更に適用できる。

最後に次数付き表示の降下は
命題\ref{prop:h3-presentation-canonical-degree-descent}を使う。
coherent tailsの平坦降下では
命題\ref{prop:h3-presented-qgr-fpqc}の強noetherian性が上で確認できたので、
構造対象とshiftも含めて有効降下する。
\end{proof}
